%==============================================================
\section{Συμπεράσματα και Μελλοντική Εργασία}
%==============================================================

Η παρούσα ενότητα ολοκληρώνει την εκτεταμένη περίληψη της διπλωματικής εργασίας. Πραγματοποιείται μια συνθετική αποτίμηση των αρχιτεκτονικών συνεισφορών του συστήματος AILS-NTUA, δίνονται ρητές απαντήσεις στα θεμελιώδη ερευνητικά ερωτήματα που καθοδήγησαν την έρευνα, αναλύονται οι εγγενείς περιορισμοί της παρούσας υλοποίησης και χαράσσονται οι μελλοντικές ερευνητικές κατευθύνσεις για τον χώρο των πολύστροφων συστημάτων RAG.

%--------------------------------------------------------------
\subsection{Σύνοψη Συνεισφορών}
%--------------------------------------------------------------

Η παρούσα εργασία διερεύνησε συστηματικά το πρόβλημα της πολύστροφης Ανάκτησης-Επαυξημένης Παραγωγής (Multi-Turn RAG) υπό το πρίσμα της σταθερότητας ανάκτησης συμφραζομένων (context-aware retrieval stability) και της αξιόπιστης πραγματολογικής αγκύρωσης αποδείξεων. Σε αντίθεση με τις συμβατικές προσεγγίσεις της βιβλιογραφίας που περιορίζονται στη βελτιστοποίηση μεμονωμένων σταδίων, αναπτύχθηκε μια ενοποιημένη αρχιτεκτονική που διαχωρίζει λειτουργικά τη διατήρηση του διαλογικού ιστορικού, την ανακάλυψη αποδείξεων και την παραγωγή απάντησης υπό περιορισμούς.

Οι κύριες ερευνητικές και τεχνολογικές συνεισφορές της εργασίας συνοψίζονται ως εξής:
\begin{itemize}
    \item \textbf{Αναδιατύπωση ως Πρωτεύον Σήμα Ανάκτησης (Query Diversity):} Στο development set και με βάση τις επίσημες κρίσεις συνάφειας του benchmark, η εισαγωγή στοχευμένης ποικιλότητας ερωτημάτων μέσω πέντε συμπληρωματικών LLM προτροπών πάνω σε έναν ενιαίο, corpus-ευθυγραμμισμένο ανακτητή (ELSER v1) υπερτερεί της σύνθεσης ετερογενών μοντέλων ανάκτησης (retriever ensembles) (Ενότητα 1.5.2).
    \item \textbf{Αλγόριθμος Ιεραρχικής Σύντηξης Κατατάξεων (Nested RRF):} Σχεδιάστηκε μια στρατηγική σύντηξης δύο επιπέδων, η οποία συγκεντρώνει τις διερευνητικές στρατηγικές (HyDE, Chain-of-Thought, Anchor-Keyword) σε μια ενδιάμεση κατάταξη πριν τη συνδυάσει με τις συντηρητικές. Η πλήρης διάταξη ανάκτησης περιορίζει την πτώση της απόδοσης μετά την πρώτη στροφή από $51\%$ σε $39\%$, χωρίς όμως να έχει απομονωθεί η συνεισφορά της ιεραρχικής σύντηξης αυτής καθαυτής.
    \item \textbf{Διύλιση Πλαισίου μέσω Evidence Span Extraction:} Η εισαγωγή ενός ενδιάμεσου σταδίου verbatim εξαγωγής αποδεικτικών προτάσεων πριν από την παραγωγή περιορίζει την επίδραση των σημασιολογικών παραπλανητών (distractors) και αποτελεί τη μεμονωμένα μεγαλύτερη συνεισφορά στον αγωγό παραγωγής ($+0.034$ HM στο development set).
    \item \textbf{Agentic Επιλογή Απάντησης Δύο Κριτών:} Ένα σχήμα εξειδικευμένων LLM κριτών (Technical, User Satisfaction) σε συνδυασμό με όρο Διαμόρφωσης Εκχύλισης (Extractiveness Shaping) στοχεύει στην εξισορρόπηση της πραγματολογικής πιστότητας με τη φυσικότητα της συνομιλιακής ροής.
    \item \textbf{Ανάλυση του Task C:} Η πτώση Task B$\rightarrow$C στο test set ($HM: 0.7698 \rightarrow 0.5409$) υποδεικνύει ως καθοριστικό παράγοντα τη βαθμονόμηση της πύλης απαντησιμότητας υπό σχεδόν τριπλάσιο ποσοστό μη απαντήσιμων ερωτημάτων, ενώ τα σφάλματα ανάκτησης στις κορυφαίες θέσεις συμβάλλουν σημαντικά· στο development set, όπου οι μη απαντήσιμες στροφές είναι σπάνιες, τα σφάλματα ανάκτησης στις κορυφαίες θέσεις εξηγούν το $72\%$ των αποτυχιών που εξετάστηκαν χειροκίνητα· καμία διαμόρφωση πύλης δεν ξεπερνά το $25\%$ F1 στο πλήρες development set (Πίνακας~\ref{tab:answerability-greek}).
\end{itemize}

%--------------------------------------------------------------
\subsection{Απαντήσεις στα Ερευνητικά Ερωτήματα}
%--------------------------------------------------------------

Τα τέσσερα κεντρικά ερευνητικά ερωτήματα (Research Questions) που διατυπώθηκαν στην Ενότητα 1.3 απαντώνται βάσει των πειραματικών δεδομένων ως εξής:
\begin{itemize}
    \item[\textbf{RQ1:}] \textit{Μπορεί η εισαγωγή στοχευμένης ποικιλότητας ερωτημάτων να υποκαταστήσει τη σύνθεση ετερογενών ανακτητών;} \\
    \textbf{Απάντηση:} Ναι, με βάση τις επίσημες κρίσεις συνάφειας. Η ποικιλότητα ερωτημάτων σε έναν ενιαίο ανακτητή υπερτερεί στο development set και αντιστοιχεί στην 1η θέση στο Task A ($nDCG@5: 0.5776$). Αντίθετα, η προσθήκη εναλλακτικών ανακτητών μειώνει το $Recall@5$, καθώς τα χρυσά έγγραφα που φέρνουν εμφανίζονται σε μέση θέση από $27{,}5$ έως $54{,}2$ μετά τη σύντηξη (Πίνακας~\ref{tab:ensemble-paradox-greek}), κάτω δηλαδή από το όριο top-$10$. Μέρος του φαινομένου μπορεί να οφείλεται στη μεροληψία των κρίσεων συνάφειας υπέρ του ELSER.
    \item[\textbf{RQ2:}] \textit{Σε ποιο βαθμό η εξαγωγή verbatim αποδεικτικών εκτάσεων συμβάλλει στον περιορισμό των ψευδαισθήσεων;} \\
    \textbf{Απάντηση:} Στο development set, η εξαγωγή spans αποτελεί τη μεγαλύτερη μεμονωμένη βελτίωση στον αγωγό παραγωγής ($+0.034$ Αρμονικό Μέσο), μαζί με αύξηση $+0.041$ στη μετρική πιστότητας χωρίς αναφορά ($RL_F$). Επειδή μια μεταβολή σε μία μόνο συνιστώσα αυξάνει τον αρμονικό μέσο τριών μετρικών μόνο κατά ένα κλάσμα της, το κέρδος $+0.034$ στον HM δείχνει ότι βελτιώθηκαν και οι μετρικές με αναφορά· η ευφράδεια δεν μετρήθηκε χωριστά. Ο περιορισμός της εισόδου του generator σε επαληθεύσιμες προτάσεις περιορίζει τον χώρο στον οποίο μπορούν να εισαχθούν αστήρικτοι ισχυρισμοί· το υψηλό $RL_F$ του test set ($0.897$) είναι συνεπές με αυτό το εύρημα.
    \item[\textbf{RQ3:}] \textit{Ποια είναι η επίδραση της κατανομικής μετατόπισης των unanswerable ερωτημάτων στην ευαισθησία των πυλών;} \\
    \textbf{Απάντηση:} Ουσιαστική. Η βαθμονόμηση της πύλης στο development set (κατώφλι $\tau{=}0.7$, ποσοστό unanswerable $6.5\%$) είναι πιθανότατα ακατάλληλη για το test set, όπου το ποσοστό σχεδόν τριπλασιάζεται ($19.1\%$), και η πτώση Task B$\rightarrow$C προέρχεται κυρίως από τη μετρική $RB_{alg}$ ($0.633 \rightarrow 0.400$). Τα στοιχεία αυτά αναδεικνύουν τη βαθμονόμηση της απαντησιμότητας ως καθοριστικό παράγοντα της end-to-end απόδοσης υπό την κατανομή του test set, αν και συμβάλλουν και τα σφάλματα ανάκτησης στις κορυφαίες θέσεις· οι δύο επιδράσεις δεν μπορούν να διαχωριστούν χωρίς ανάλυση των σκορ του test set ανά κατηγορία. Το σχήμα πολλαπλών κριτών βελτιώνει τον μονό κριτή στο development set, αλλά αφήνει ισχυρό acceptance bias.
    \item[\textbf{RQ4:}] \textit{Μπορεί η επαναβαθμολόγηση μέσω ΜΓΜ (LLM-based reranking) να προσφέρει όφελος έναντι εξειδικευμένων cross-encoder μοντέλων όταν η ανάκληση του πρώτου σταδίου είναι κοντά στον κορεσμό;} \\
    \textbf{Απάντηση:} Όχι ουσιαστικά. Καμία από τις τρεις μορφές LLM reranking (pointwise, listwise, generation-based) δεν βελτίωσε τη σταθμισμένη σύντηξη ανακτητή και cross-encoder όταν το $Recall@100$ ξεπερνά το $0.90$· ως βοηθητικό σήμα με μικρό βάρος δίνουν το πολύ οριακά κέρδη (έως $+2$ μονάδες $Recall@5$ ανά θεματικό πεδίο), που δεν δικαιολογούν το πρόσθετο κόστος.
\end{itemize}

%--------------------------------------------------------------
\subsection{Περιορισμοί του Συστήματος}
%--------------------------------------------------------------

Παρά τις υψηλές επιδόσεις στα Tasks A και B, το σύστημα και η αξιολόγησή του έχουν συγκεκριμένους περιορισμούς:
\begin{enumerate}
    \item \textbf{Ρύθμιση και αναφορά στο ίδιο σύνολο:} Όλες οι υπερπαράμετροι ρυθμίστηκαν στο development set, στο οποίο αναφέρονται και όλα τα ablations, χωρίς ξεχωριστό τμήμα ελέγχου ή cross-validation ανά συνομιλία· τα αποτελέσματα του development set είναι επομένως πιθανώς αισιόδοξα. Κάθε διαμόρφωση εκτελέστηκε μία φορά, οπότε μικρές διαφορές (της τάξης του $0.01$ HM ή λιγότερο) δεν πρέπει να υπερερμηνεύονται.
    \item \textbf{Στατικότητα Κατωφλίων Απόφασης:} Η Πύλη Απαντησιμότητας βασίζεται σε σταθερό κατώφλι εμπιστοσύνης ($\tau=0.7$), βαθμονομημένο αποκλειστικά στο development set. Το κατώφλι αυτό, μαζί με τη ζώνη εκχύλισης ($[0.28, 0.38]$) και τα βάρη του Nested RRF ανά corpus, δεν επανα-βαθμονομήθηκαν για το test set. Η αναντιστοιχία αυτή είναι πιθανότατα σημαντικός παράγοντας του χάσματος Task B$\rightarrow$C, μαζί με τα σφάλματα ανάκτησης στις κορυφαίες θέσεις· το αν ένα άλλο κατώφλι θα απέδιδε καλύτερα στο test set δεν μπορεί να επαληθευτεί με μία μόνο υποβολή.
    \item \textbf{Ιστορικό αναφοράς στην αξιολόγηση:} Το MTRAGEval παρέχει σε κάθε στροφή το ιστορικό αναφοράς, οπότε η αξιολόγηση δεν ελέγχει την ανθεκτικότητα του συστήματος στα δικά του προηγούμενα σφάλματα, όπως θα συνέβαινε σε έναν πραγματικό βοηθό.
    \item \textbf{Μεροληψία των κρίσεων συνάφειας:} Οι κρίσεις συνάφειας του MTRAG παρήχθησαν κυρίως από αποσπάσματα του ELSER, γεγονός που μπορεί να ευνοεί το ELSER έναντι άλλων ανακτητών και να συμβάλλει στο παράδοξο σύνθεσης ανακτητών.
    \item \textbf{Εξάρτηση από Κλειστά API (Proprietary Models):} Η αρχιτεκτονική βασίζεται στο GPT-4o και στο DeepSeek-V3.2, γεγονός που αυξάνει το λειτουργικό κόστος και μειώνει την αναπαραγωγιμότητα, λόγω της έλλειψης πρόσβασης στα βάρη των μοντέλων. Κανένα στάδιο του αγωγού δεν έχει υποστεί fine-tuning στα δεδομένα MTRAG· όλα τα κέρδη προέρχονται από μηχανική προτροπών (prompt engineering) και αρχιτεκτονικές επιλογές.
    \item \textbf{Υπολογιστική Πολυπλοκότητα Pipeline:} Η πολυσταδιακή φύση του αγωγού (5 αναδιατυπώσεις, cross-encoder reranking, πύλη απαντησιμότητας τριών κριτών) αυξάνει τη δομική πολυπλοκότητα και τον χρόνο απόκρισης, ο οποίος δεν μετρήθηκε, και απαιτεί πρόσθετη μηχανική ενορχήστρωσης (orchestration) για χρήση σε κλίμακα.
\end{enumerate}

%--------------------------------------------------------------
\subsection{Μελλοντικές Ερευνητικές Κατευθύνσεις}
%--------------------------------------------------------------

Η παρούσα διπλωματική εργασία ανοίγει τις ακόλουθες ερευνητικές κατευθύνσεις:
\begin{itemize}
    \item \textbf{Προσαρμοστική Βαθμονόμηση (Adaptive Thresholding):} Δυναμικοί μηχανισμοί απαντησιμότητας, όπου το κατώφλι $\tau$ μεταβάλλεται βάσει της εκτίμησης της εντροπίας των ανακτηθέντων εγγράφων και της ανίχνευσης μετατοπίσεων κατανομής, αντί να παραμένει στατικό. Επειδή η βαθμονόμηση της απαντησιμότητας είναι ο καθοριστικός παράγοντας υπό την κατανομή του test set, πρόκειται για κατεύθυνση υψηλής αναμενόμενης αξίας· η ανάλυση των επίσημων σκορ του test set ανά κατηγορία απαντησιμότητας θα βοηθούσε να ποσοτικοποιηθεί το αναμενόμενο κέρδος.
    \item \textbf{Ισχυρότερες γραμμές βάσης για την πύλη:} Ταξινομητές επάρκειας βασισμένοι σε NLI, ταξινομητές πάνω στα σκορ ανάκτησης και κατώφλια βαθμονομημένα ανά θεματικό πεδίο, που δεν αξιολογήθηκαν στην παρούσα εργασία.
    \item \textbf{Ανθεκτική επίλυση εξαρτήσεων μεταξύ στροφών:} Η δεύτερη μεγαλύτερη κατηγορία σφαλμάτων είναι η ανεπίλυτη εξάρτηση από προηγούμενες στροφές ($34\%$). Καλύτερη επίλυση αναφορών και ελλείψεων ---π.χ. με ρητή παρακολούθηση οντοτήτων--- στοχεύει άμεσα σε αυτήν. Σε πραγματικά συστήματα, όπου το ιστορικό περιέχει τις απαντήσεις του ίδιου του συστήματος, μηχανισμοί ελέγχου και αναίρεσης προηγούμενων βημάτων (backtracking) θα μπορούσαν επιπλέον να περιορίσουν το error cascade.
    \item \textbf{Fine-Tuning Ανοιχτών Μοντέλων Συλλογιστικής (Open-Weight Reasoning):} Εκπαίδευση εξειδικευμένων open-weight μοντέλων (π.χ. Llama-3-Reasoning, DeepSeek-R1) μέσω απόσταξης γνώσης (knowledge distillation) από το προτεινόμενο pipeline, με στόχο τη μείωση του κόστους και την άρση της εξάρτησης από κλειστά API.
    \item \textbf{Ενισχυτική Μάθηση (RL) για Query Formulation:} Εφαρμογή αλγορίθμων ενισχυτικής μάθησης (π.χ. PPO) για την εκπαίδευση του query rewriter, με συνάρτηση ανταμοιβής (reward function) απευθείας το nDCG@5 του ανακτητή.
\end{itemize}

%--------------------------------------------------------------
\subsection{Τελικό Συμπέρασμα}
%--------------------------------------------------------------

Συμπερασματικά, η παρούσα εργασία υποδεικνύει ότι η επιτυχής υλοποίηση συστημάτων Πολύστροφου RAG δεν εξαρτάται αποκλειστικά από τη χρήση ισχυρότερων γλωσσικών μοντέλων. Η μετατόπιση της έμφασης από την απλή κλιμάκωση (scaling) προς τον προσεκτικό αρχιτεκτονικό σχεδιασμό, τη \textit{σταθερότητα της ανάκτησης} και την \textit{αυστηρή βαθμονόμηση των αποφάσεων} αποτελεί μια υποσχόμενη κατεύθυνση για συστήματα υψηλής αξιοπιστίας και διαφάνειας. Τα ευρήματα της εργασίας στηρίζονται στις αναλύσεις του development set και στα επίσημα αποτελέσματα, με την 1η και τη 2η θέση στα Tasks A και B· η υστέρηση στο Task C αναδεικνύει τη \textit{βαθμονόμηση της απαντησιμότητας}, δηλαδή την απόφαση απάντησης ή αποχής υπό μετατόπιση κατανομής, ως το καθοριστικό ανοιχτό πρόβλημα, με την ακρίβεια κορυφής της ανάκτησης να παραμένει σημαντικός δευτερεύων περιορισμός.
